\section{Related Work}
\label{sec:related}

We situate HADL at the intersection of five lines of prior work, spanning LLM agent orchestration frameworks, programming language approaches to LLM interaction, gradual type systems, metacircular interpretation, and policy languages for authorization.

\subsection{LLM Agent Frameworks}

The dominant approach to building agentic systems today is through orchestration libraries. LangGraph~\cite{langgraph2024} models workflows as state machines with durable execution and human-in-the-loop interrupts, while LangChain~\cite{langchain2022} provides composable abstractions for chaining LLM calls with tool integrations. AutoGen~\cite{wu2024autogen} introduces multi-agent conversation protocols in which agents coordinate through message passing, and MetaGPT~\cite{hong2023metagpt} assigns specialized roles to agents that collaborate through structured outputs. These systems have demonstrated practical value in production settings and offer rich integrations with LLM providers and external services. Beyond agent-specific frameworks, durable execution platforms such as Temporal~\cite{temporal2024} provide infrastructure for long-running, fault-tolerant workflows with state persistence and replay. These systems offer robust operational semantics but are general-purpose workflow engines without type systems tailored to LLM interaction or agent-specific primitives. All of these approaches share a common limitation from a programming language perspective. None provides a static type system, so mismatches between expected and actual tool outputs are caught only at runtime. None offers compile-time error reporting with source locations. Control flow is expressed through graph construction or callback registration rather than through imperative constructs such as loops and conditionals, and authorization is handled externally if at all. HADL addresses these limitations by treating agentic plans as typed programs rather than as library configurations, providing gradual type checking, compile-time diagnostics, and policy enforcement as language-level features.

\subsection{Programming Language Approaches to LLM Interaction}

A smaller but growing body of work applies programming language ideas directly to LLM interaction. LMQL~\cite{beurer2023prompting}, published at PLDI 2023, is the closest precedent in the PL literature. It introduces a query language for LLMs with scripted prompting, constrained decoding, and type-like constraints on outputs, demonstrating that PL techniques can improve the reliability of LLM interactions. However, LMQL targets single-turn prompt scripting rather than multi-turn agent orchestration, and does not support mutable state, unbounded loops, or workflows as composable typed functions. DSPy~\cite{khattab2023dspy} takes a compiler-oriented approach, treating LLM pipelines as declarative programs that can be optimized through automatic prompt tuning and few-shot selection. DSPy's compilation model is powerful for optimizing prompting strategies but operates at the pipeline level rather than the expression level, does not provide a static type system, and is not designed for the long-running stateful execution patterns that characterize agentic workflows. Daunis~\cite{daunis2025declarative} proposes a declarative DSL for enterprise agent workflows that separates specification from implementation using DAG-based pipeline definitions, achieving significant reductions in development time at PayPal. This system targets a different design point from HADL. It is declarative rather than imperative, uses DAG configuration rather than a type-checked language, and cannot express unbounded iteration or mutable state accumulation. HADL is implemented as a standalone DSL with its own PEG grammar, parser, and file format (\texttt{.hadl}), providing domain-specific syntax and dedicated static checking. HADL occupies a distinct position in this design space by combining gradual type checking with imperative control flow, mutable state, and integrated policy enforcement within a single language.

\subsection{Gradual Type Checking and Structural Subtyping}

HADL's type system builds on the theory of gradual typing introduced by Siek and Taha~\cite{siek2006gradual}, who formalized a discipline for integrating static and dynamic typing within a single language. In their framework, a special unknown type allows regions of a program to defer type checking to runtime, with casts inserted at the boundary between statically and dynamically typed code. Siek et al.~\cite{siek2015refined} subsequently refined the criteria for gradual typing, introducing the \emph{gradual guarantee}, which states that adding or removing type annotations preserves well-typedness and does not change program behavior beyond introducing or removing cast errors. Tobin-Hochstadt and Felleisen~\cite{tobinhochstadt2006interlanguage} developed interlanguage migration techniques for gradually introducing types into untyped Scheme modules, and their subsequent work on Typed Racket~\cite{tobinhochstadt2008design} provided the first practical implementation of gradual typing with sound interoperation between typed and untyped code. Vitousek et al.~\cite{vitousek2017big} addressed the challenge of maintaining type soundness in open-world settings where gradually typed code interacts with untyped external libraries.

HADL adapts gradual typing to a new domain in which the source of dynamic types is not programmer choice but LLM execution. In standard gradual typing, the programmer decides which regions of code are statically or dynamically typed. In HADL, the boundary arises from the nature of LLM tool calls, which can generate both values and types that do not exist at compile time. HADL's staged type checking mechanism extends the gradual approach by re-invoking the type checker at runtime whenever an LLM call materializes a new type, substituting the concrete type into the program and verifying the downstream code before execution resumes. This re-verifies downstream variable bindings and subtyping checks against the materialized schema before dependent code proceeds. When re-checking fails, a provenance-driven self-healing mechanism re-invokes the causal tool call with the type error as feedback. Field-level property access on schema-typed values is performed by the embedded JavaScript engine at runtime, operating on the underlying JSON values. This staged approach provides broader type coverage than standard cast-based enforcement for this specific domain, as it re-checks entire program regions rather than individual boundary crossings.

For structural records, HADL's runtime schema system follows Cardelli's~\cite{cardelli1984semantics} formalization of width and depth subtyping and R\'{e}my's~\cite{remy1989type} treatment of records with field-order independence. The schema type, which represents JSON Schema objects at runtime, implements a subtyping partial order in which an object with more properties is a subtype of one with fewer (width subtyping) and nested property types are compared covariantly (depth subtyping). Pierce~\cite{pierce2002types} provides the standard reference for the subtyping metatheory. Cardelli and Wegner~\cite{cardelli1985understanding} supply the broader conceptual framework situating structural subtyping within the landscape of polymorphism.

\subsection{Metacircular Interpretation}

The concept of metacircular interpretation, in which a language is powerful enough to express its own interpreter, originates with McCarthy's~\cite{mccarthy1960recursive} Lisp evaluator and was made rigorous by Reynolds~\cite{reynolds1972definitional}, who showed that the evaluation strategy of the meta-language determines the semantics of the object language. Landin~\cite{landin1964mechanical} demonstrated that expressions could be modeled and evaluated mechanically through $\lambda$-notation, establishing a direct precursor to definitional interpreters. Abelson and Sussman~\cite{abelson2022structure} made the metacircular evaluator central to computer science education, and Friedman and Wand~\cite{friedman2001essentials} provided a systematic treatment of interpreter-based semantics with progressive extensions for state, continuations, and effects. Amin and Rompf~\cite{amin2017collapsing} showed that towers of interpreters can be collapsed through staging, connecting metacircular interpretation to compilation. HADL draws on this tradition in two ways. First, its architect pattern uses \texttt{agent} calls to generate sub-plans expressed in HADL syntax, demonstrating that the language's own syntax can appear as the output of LLM invocations. Second, HADL provides a first-class \texttt{eval} expression that deserializes a workflow value (a JSON-serialized AST) and executes it within the current interpreter context. The \texttt{eval} expression is typed, so the type checker resolves the workflow's arrow type and checks the supplied arguments against it, or defers to gradual typing when the workflow value has type \texttt{Schema}. At runtime, \texttt{eval} deserializes the JSON value into a \texttt{Program} AST, binds the arguments to the workflow's parameters, and evaluates the body using the existing tree-walking interpreter. This provides a concrete realization of the metacircular capacity, where an LLM can generate a typed HADL workflow as structured output and the outer workflow can invoke it via \texttt{eval} with full type checking at the call site.

\subsection{Policy Languages and Authorization}

HADL integrates Cedar~\cite{cutler2024cedar}, a policy language designed to be simultaneously expressive, performant, safe, and analyzable. Cedar externalizes authorization logic from application code into standalone policies evaluated by a dedicated engine, supporting role-based and attribute-based access control through (principal, action, resource) triplets. Cedar's design has been formally verified in Lean and validated through differential testing against a Dafny model, providing high-assurance correctness guarantees. HADL's contribution relative to Cedar is one of integration rather than language design. Where Cedar provides the policy language and evaluation engine, HADL embeds Cedar's authorization model directly into the runtime evaluation loop of an agentic workflow language. Before every \texttt{tool} and \texttt{agent} invocation, the interpreter constructs an authorization triplet identifying the invoking principal, the requested action, and the target resource, and validates it against a Cedar policy store. This integration makes authorization a property of the language's execution semantics rather than an external middleware concern, addressing the coarse-grained access control problem identified in Section~\ref{sec:intro}.